\documentclass[../main]{subfiles}
\ifSubfilesClassLoaded{\setcounter{chapter}{10}}{}
\begin{document}

\chapter{Representasi Fungsi Logika ke Rangkaian Gerbang Logika Dasar}

\begin{subcpmk}
  \item \textbf{Sub-CPMK 3.3:} Mahasiswa mampu mengonversi dan merepresentasikan fungsi logika yang telah disederhanakan ke dalam bentuk rangkaian gerbang logika dasar pembentuk sistem komputer.
\end{subcpmk}

% ============================================================
% MATERI POKOK
% ============================================================
\section{Gerbang Logika Dasar: AND, OR, NOT}

Gerbang logika (\textit{logic gate}) merupakan komponen fundamental dalam arsitektur komputer dan sistem digital. Setiap gerbang logika mengimplementasikan satu operasi aljabar Boolean pada sinyal-sinyal biner, menerima satu atau lebih masukan dan menghasilkan tepat satu keluaran. Pada tingkat fisik, gerbang logika direalisasikan menggunakan transistor semikonduktor, namun pada tingkat konseptual, kita merepresentasikannya sebagai simbol grafis standar yang mengikuti konvensi IEEE/ANSI \cite{allaboutcircuits_gates}.

\subsection{Gerbang AND}

\begin{definisi}[Gerbang AND]
Gerbang AND adalah gerbang logika biner yang menghasilkan keluaran bernilai $1$ \textbf{jika dan hanya jika seluruh} masukannya bernilai $1$. Jika salah satu atau lebih masukan bernilai $0$, maka keluaran bernilai $0$. Operasi ini berkorespondensi dengan perkalian Boolean ($\cdot$) dan konjungsi ($\wedge$) pada logika proposisi.
\end{definisi}

Gerbang AND dapat memiliki dua atau lebih masukan. Untuk gerbang AND dengan $n$ masukan, keluaran bernilai $1$ hanya jika seluruh $n$ masukan bernilai $1$. Tabel kebenaran gerbang AND 2-masukan ditampilkan pada Tabel~\ref{tab:and-gate}.

\begin{table}[H]
\centering
\renewcommand{\arraystretch}{1.3}
\begin{tabular}{|c|c||c|}
\hline
$A$ & $B$ & $F = A \cdot B$ \\ \hline
0 & 0 & 0 \\ \hline
0 & 1 & 0 \\ \hline
1 & 0 & 0 \\ \hline
1 & 1 & 1 \\ \hline
\end{tabular}
\caption{Tabel Kebenaran Gerbang AND}
\label{tab:and-gate}
\end{table}

\subsection{Gerbang OR}

\begin{definisi}[Gerbang OR]
Gerbang OR adalah gerbang logika biner yang menghasilkan keluaran bernilai $1$ \textbf{jika sekurang-kurangnya satu} masukannya bernilai $1$. Keluaran bernilai $0$ hanya jika seluruh masukan bernilai $0$. Operasi ini berkorespondensi dengan penjumlahan Boolean ($+$) dan disjungsi ($\vee$) pada logika proposisi.
\end{definisi}

Analoginya, gerbang OR bekerja seperti rangkaian listrik paralel: jika salah satu sakelar ditutup (bernilai 1), arus dapat mengalir (keluaran 1). Gerbang OR juga dapat memiliki lebih dari dua masukan, di mana keluaran bernilai $1$ jika setidaknya satu dari $n$ masukan bernilai $1$.

\subsection{Gerbang NOT (Inverter)}

\begin{definisi}[Gerbang NOT]
Gerbang NOT (inverter) adalah gerbang logika uner yang memiliki tepat satu masukan dan menghasilkan keluaran yang merupakan \textbf{komplemen} (negasi) dari masukan tersebut. Jika masukan 0, keluaran 1; jika masukan 1, keluaran 0. Operasi ini berkorespondensi dengan komplemen ($'$) dan negasi ($\neg$) pada logika proposisi.
\end{definisi}

Gerbang NOT adalah satu-satunya gerbang dengan tepat satu masukan. Pada simbol standar, gerbang NOT direpresentasikan sebagai segitiga dengan lingkaran kecil (\textit{bubble}) pada keluarannya. Lingkaran ini menandakan operasi inversi dan juga digunakan pada simbol gerbang turunan seperti NAND dan NOR.

\subsection{Simbol Standar dan Representasi}

Gambar~\ref{fig:gerbang-dasar-11} menampilkan simbol grafis standar IEEE/ANSI untuk ketiga gerbang logika dasar. Simbol-simbol ini digunakan secara universal dalam diagram rangkaian digital.

\begin{figure}[H]
\centering
\begin{tikzpicture}[circuit logic US, huge circuit symbols,
                    every circuit symbol/.style={fill=blue!10, draw=blue!50!black, thick}]

  % Gerbang AND
  \node[and gate] (and1) at (0,0) {};
  \draw (and1.input 1) -- ++(-0.5,0) node[left] {$A$};
  \draw (and1.input 2) -- ++(-0.5,0) node[left] {$B$};
  \draw (and1.output) -- ++(0.5,0) node[right] {$A \cdot B$};
  \node[below=0.5cm of and1] {AND};

  % Gerbang OR
  \node[or gate] (or1) at (4.5,0) {};
  \draw (or1.input 1) -- ++(-0.5,0) node[left] {$A$};
  \draw (or1.input 2) -- ++(-0.5,0) node[left] {$B$};
  \draw (or1.output) -- ++(0.5,0) node[right] {$A + B$};
  \node[below=0.5cm of or1] {OR};

  % Gerbang NOT
  \node[not gate] (not1) at (9,0) {};
  \draw (not1.input) -- ++(-0.5,0) node[left] {$A$};
  \draw (not1.output) -- ++(0.5,0) node[right] {$A'$};
  \node[below=0.5cm of not1] {NOT};

\end{tikzpicture}
\caption{Simbol Standar IEEE/ANSI Gerbang Logika Dasar}
\label{fig:gerbang-dasar-11}
\end{figure}

Dalam praktik perancangan rangkaian, ketiga gerbang dasar ini sudah mencukupi untuk mengimplementasikan fungsi Boolean apa pun, karena himpunan $\{AND, OR, NOT\}$ bersifat lengkap secara fungsional (\textit{functionally complete}). Setiap fungsi Boolean dapat ditulis dalam bentuk Sum of Products (SOP) yang hanya menggunakan operator AND, OR, dan NOT \cite{rosen2019}. Namun, dalam praktik produksi semikonduktor, gerbang turunan dan universal lebih sering digunakan karena alasan efisiensi fabrikasi, sebagaimana akan dibahas pada section berikutnya.

\section{Gerbang Logika Turunan: NAND, NOR, XOR, XNOR}

Selain tiga gerbang dasar, terdapat empat gerbang logika turunan yang penting dalam perancangan digital. Gerbang-gerbang ini disebut ``turunan'' karena fungsinya dapat diekspresikan menggunakan kombinasi gerbang AND, OR, dan NOT. Namun, dalam implementasi fisik, gerbang turunan justru seringkali memerlukan lebih sedikit transistor dibandingkan kombinasi gerbang dasar yang ekuivalennya, menjadikannya lebih efisien secara perangkat keras \cite{allaboutcircuits_gates}.

\subsection{Gerbang NAND (NOT-AND)}

\begin{definisi}[Gerbang NAND]
Gerbang NAND menghasilkan keluaran yang merupakan \textbf{komplemen} dari gerbang AND. Keluaran bernilai $0$ hanya jika seluruh masukan bernilai $1$, dan bernilai $1$ pada kondisi lainnya. Ekspresi Boolean: $F = (A \cdot B)' = A' + B'$ (berdasarkan hukum De Morgan).
\end{definisi}

Gerbang NAND memiliki signifikansi khusus karena merupakan salah satu dari dua gerbang universal yang akan dibahas pada section berikutnya. Pada level transistor, gerbang NAND merupakan gerbang yang paling natural dalam teknologi CMOS (\textit{Complementary Metal-Oxide-Semiconductor}), karena implementasinya hanya memerlukan empat transistor, lebih sedikit dibandingkan implementasi AND yang memerlukan enam transistor.

\subsection{Gerbang NOR (NOT-OR)}

\begin{definisi}[Gerbang NOR]
Gerbang NOR menghasilkan keluaran yang merupakan \textbf{komplemen} dari gerbang OR. Keluaran bernilai $1$ hanya jika seluruh masukan bernilai $0$, dan bernilai $0$ jika setidaknya satu masukan bernilai $1$. Ekspresi Boolean: $F = (A + B)' = A' \cdot B'$ (berdasarkan hukum De Morgan).
\end{definisi}

Gerbang NOR juga merupakan gerbang universal. Secara historis, komputer panduan Apollo (\textit{Apollo Guidance Computer}) yang digunakan NASA untuk misi pendaratan di bulan tahun 1969 dibangun seluruhnya dari sekitar 5.600 gerbang NOR \cite{electronics_tutorials_gates}.

\subsection{Gerbang XOR (Exclusive-OR)}

\begin{definisi}[Gerbang XOR]
Gerbang XOR (\textit{Exclusive-OR}) menghasilkan keluaran bernilai $1$ jika dan hanya jika \textbf{jumlah masukan bernilai $1$ bersifat ganjil}. Untuk gerbang 2-masukan, keluaran bernilai $1$ hanya jika kedua masukan memiliki nilai yang \textbf{berbeda}. Ekspresi Boolean: $F = A \oplus B = A'B + AB'$.
\end{definisi}

Gerbang XOR memiliki aplikasi yang luas dalam sistem digital:
\begin{itemize}
    \item \textbf{Penjumlah biner}: XOR menghasilkan bit jumlah (sum) pada rangkaian half-adder dan full-adder.
    \item \textbf{Detektor paritas}: XOR dari beberapa bit menghitung paritas (genap/ganjil) dari data.
    \item \textbf{Kriptografi}: Operasi XOR digunakan secara ekstensif dalam algoritma enkripsi seperti \textit{stream cipher} dan \textit{one-time pad}.
    \item \textbf{Komparator}: XOR bernilai 1 jika dua bit berbeda, sehingga dapat digunakan untuk mendeteksi perbedaan.
\end{itemize}

\subsection{Gerbang XNOR (Exclusive-NOR)}

\begin{definisi}[Gerbang XNOR]
Gerbang XNOR adalah komplemen dari gerbang XOR. Keluaran bernilai $1$ jika dan hanya jika kedua masukan memiliki nilai yang \textbf{sama} (keduanya 0 atau keduanya 1). Ekspresi Boolean: $F = (A \oplus B)' = AB + A'B' = A \odot B$.
\end{definisi}

Gerbang XNOR berfungsi sebagai detektor kesamaan (\textit{equality detector}) dan sering digunakan dalam rangkaian komparator digital, di mana beberapa pasang bit dibandingkan secara paralel untuk menentukan apakah dua bilangan biner bernilai sama.

\subsection{Simbol dan Tabel Kebenaran Perbandingan}

Gambar~\ref{fig:gerbang-turunan-11} menampilkan simbol grafis keempat gerbang turunan.

\begin{figure}[H]
\centering
\begin{tikzpicture}[circuit logic US, huge circuit symbols,
                    every circuit symbol/.style={thick}]

  % Gerbang NAND
  \node[nand gate, fill=cyan!10, draw=cyan!50!black] (nand1) at (0,0) {};
  \draw (nand1.input 1) -- ++(-0.5,0) node[left] {$A$};
  \draw (nand1.input 2) -- ++(-0.5,0) node[left] {$B$};
  \draw (nand1.output) -- ++(0.5,0) node[right] {$(AB)'$};
  \node[below=0.5cm of nand1] {NAND};

  % Gerbang NOR
  \node[nor gate, fill=cyan!10, draw=cyan!50!black] (nor1) at (4.5,0) {};
  \draw (nor1.input 1) -- ++(-0.5,0) node[left] {$A$};
  \draw (nor1.input 2) -- ++(-0.5,0) node[left] {$B$};
  \draw (nor1.output) -- ++(0.5,0) node[right] {$(A+B)'$};
  \node[below=0.5cm of nor1] {NOR};

  % Gerbang XOR
  \node[xor gate, fill=orange!10, draw=orange!50!black] (xor1) at (0,-3) {};
  \draw (xor1.input 1) -- ++(-0.5,0) node[left] {$A$};
  \draw (xor1.input 2) -- ++(-0.5,0) node[left] {$B$};
  \draw (xor1.output) -- ++(0.5,0) node[right] {$A \oplus B$};
  \node[below=0.5cm of xor1] {XOR};

  % Gerbang XNOR
  \node[xnor gate, fill=orange!10, draw=orange!50!black] (xnor1) at (4.5,-3) {};
  \draw (xnor1.input 1) -- ++(-0.5,0) node[left] {$A$};
  \draw (xnor1.input 2) -- ++(-0.5,0) node[left] {$B$};
  \draw (xnor1.output) -- ++(0.5,0) node[right] {$A \odot B$};
  \node[below=0.5cm of xnor1] {XNOR};

\end{tikzpicture}
\caption{Simbol Gerbang Logika Turunan}
\label{fig:gerbang-turunan-11}
\end{figure}

Tabel~\ref{tab:perbandingan-gerbang-11} menyajikan perbandingan lengkap tabel kebenaran seluruh gerbang logika 2-masukan.

\begin{table}[H]
\centering
\small
\renewcommand{\arraystretch}{1.3}
\begin{tabular}{|c|c||c|c|c|c|c|c|}
\hline
$A$ & $B$ & \textbf{AND} & \textbf{OR} & \textbf{NAND} & \textbf{NOR} & \textbf{XOR} & \textbf{XNOR} \\ \hline
0 & 0 & 0 & 0 & 1 & 1 & 0 & 1 \\ \hline
0 & 1 & 0 & 1 & 1 & 0 & 1 & 0 \\ \hline
1 & 0 & 0 & 1 & 1 & 0 & 1 & 0 \\ \hline
1 & 1 & 1 & 1 & 0 & 0 & 0 & 1 \\ \hline
\end{tabular}
\caption{Perbandingan Tabel Kebenaran Seluruh Gerbang Logika 2-Masukan}
\label{tab:perbandingan-gerbang-11}
\end{table}

\begin{catatan}
Perhatikan hubungan komplementer pada tabel di atas: NAND adalah komplemen AND, NOR adalah komplemen OR, dan XNOR adalah komplemen XOR. Setiap pasangan gerbang komplementer menghasilkan keluaran yang berlawanan untuk setiap kombinasi masukan.
\end{catatan}

\section{Gerbang Universal dan Kelengkapan Fungsional}

Konsep kelengkapan fungsional (\textit{functional completeness}) merupakan salah satu hasil sentral dalam teori rangkaian digital. Suatu himpunan gerbang logika dikatakan lengkap secara fungsional jika setiap fungsi Boolean dapat diimplementasikan hanya menggunakan gerbang-gerbang dalam himpunan tersebut. Secara praktis, gerbang NAND secara tunggal bersifat lengkap secara fungsional, dan demikian pula gerbang NOR. Kedua gerbang ini disebut \logic{gerbang universal} \cite{rosen2019}.

\subsection{Pembuktian Universalitas NAND}

Untuk membuktikan bahwa NAND bersifat universal, cukup ditunjukkan bahwa ketiga operasi dasar (NOT, AND, OR) dapat diimplementasikan hanya menggunakan gerbang NAND. Jika $\uparrow$ melambangkan operasi NAND, maka $(A \uparrow B) = (A \cdot B)'$.

\begin{enumerate}[label=\textbf{\arabic*.}]
  \item \textbf{NOT dari NAND}: Hubungkan kedua masukan NAND ke sinyal yang sama:
  \[ A \uparrow A = (A \cdot A)' = A' \]

  \item \textbf{AND dari NAND}: AND adalah komplemen dari NAND, sehingga dibutuhkan dua gerbang NAND:
  \[ A \cdot B = ((A \cdot B)')' = (A \uparrow B) \uparrow (A \uparrow B) \]

  \item \textbf{OR dari NAND}: Berdasarkan hukum De Morgan, $A + B = (A' \cdot B')'$:
  \[ A + B = (A' \cdot B')' = (A \uparrow A) \uparrow (B \uparrow B) \]
  Diperlukan tiga gerbang NAND: dua untuk inversi dan satu untuk operasi NAND pada hasil inversi.
\end{enumerate}

\begin{figure}[H]
\centering
\begin{tikzpicture}[circuit logic US, huge circuit symbols,
                    every circuit symbol/.style={fill=cyan!10, draw=cyan!50!black, thick},
                    scale=0.85, transform shape]

  % NOT from NAND
  \node[font=\bfseries] at (0, 1.5) {NOT dari NAND};
  \node[nand gate] (n1) at (1.5, 0) {};
  \draw (n1.input 1) -- ++(-0.5,0) coordinate (j1);
  \draw (n1.input 2) -- ++(-0.5,0) coordinate (j2);
  \draw (j1) -- ++(-0.5, 0) |- (j2);
  \draw (j1) -- ++(-0.5,0) -- ++(-0.3,0) node[left] {$A$};
  \draw (n1.output) -- ++(0.5,0) node[right] {$A'$};

  % AND from NAND
  \node[font=\bfseries] at (6, 1.5) {AND dari NAND};
  \node[nand gate] (n2) at (7, 0) {};
  \node[nand gate] (n3) at (9.5, 0) {};
  \draw (n2.input 1) -- ++(-0.5,0) node[left] {$A$};
  \draw (n2.input 2) -- ++(-0.5,0) node[left] {$B$};
  \draw (n2.output) -- (n3.input 1);
  \draw (n2.output) ++(0.3,0) |- (n3.input 2);
  \draw (n3.output) -- ++(0.5,0) node[right] {$AB$};

\end{tikzpicture}
\caption{Implementasi NOT dan AND menggunakan Gerbang NAND}
\label{fig:nand-universal}
\end{figure}

\subsection{Pembuktian Universalitas NOR}

Dengan cara yang analog, universalitas NOR dibuktikan dengan mengimplementasikan NOT, OR, dan AND hanya menggunakan gerbang NOR. Jika $\downarrow$ melambangkan operasi NOR, maka $(A \downarrow B) = (A + B)'$.

\begin{enumerate}[label=\textbf{\arabic*.}]
  \item \textbf{NOT dari NOR}: $A \downarrow A = (A + A)' = A'$
  \item \textbf{OR dari NOR}: $A + B = ((A + B)')' = (A \downarrow B) \downarrow (A \downarrow B)$
  \item \textbf{AND dari NOR}: $A \cdot B = (A' + B')' = (A \downarrow A) \downarrow (B \downarrow B)$
\end{enumerate}

\subsection{Perbandingan Gerbang yang Bersifat Universal}

Tidak semua gerbang bersifat universal. Tabel~\ref{tab:universalitas} menunjukkan bahwa dari tujuh jenis gerbang logika, hanya NAND dan NOR yang secara individual bersifat lengkap secara fungsional.

\begin{table}[H]
\centering
\renewcommand{\arraystretch}{1.3}
\begin{tabular}{|l|c|l|}
\hline
\textbf{Gerbang} & \textbf{Universal?} & \textbf{Alasan} \\ \hline
AND & Tidak & Tidak dapat menghasilkan NOT \\ \hline
OR & Tidak & Tidak dapat menghasilkan NOT \\ \hline
NOT & Tidak & Tidak dapat menghasilkan AND/OR (hanya uner) \\ \hline
\textbf{NAND} & \textbf{Ya} & Dapat menghasilkan NOT, AND, OR \\ \hline
\textbf{NOR} & \textbf{Ya} & Dapat menghasilkan NOT, AND, OR \\ \hline
XOR & Tidak & Tidak dapat menghasilkan AND \\ \hline
XNOR & Tidak & Tidak dapat menghasilkan AND \\ \hline
\end{tabular}
\caption{Status Kelengkapan Fungsional Gerbang Logika}
\label{tab:universalitas}
\end{table}

\subsection{Implikasi Praktis dalam Produksi}

Universalitas NAND dan NOR memiliki konsekuensi praktis yang besar bagi industri semikonduktor. Dalam proses fabrikasi chip (\textit{wafer fabrication}), memproduksi satu jenis gerbang secara massal jauh lebih efisien dibandingkan memproduksi beragam jenis gerbang pada satu chip. Keluarga logika TTL (\textit{Transistor-Transistor Logic}) menggunakan NAND sebagai komponen primitif, sementara keluarga logika ECL (\textit{Emitter-Coupled Logic}) menggunakan NOR \cite{electronics_tutorials_gates}.

Selain efisiensi fabrikasi, penggunaan satu jenis gerbang juga menyederhanakan proses verifikasi dan pengujian. Alat-alat EDA modern secara otomatis mengonversi desain dari bentuk AND-OR-NOT ke dalam bentuk NAND-only atau NOR-only melalui proses yang disebut \logic{technology mapping}. Proses ini merupakan bagian integral dari alur desain VLSI (\textit{Very Large Scale Integration}) yang menghasilkan chip-chip terintegrasi modern.

\section{Konversi Ekspresi Boolean ke Rangkaian Gerbang}

Proses mentranslasikan ekspresi aljabar Boolean menjadi diagram rangkaian gerbang logika merupakan kompetensi utama dalam perancangan digital. Setiap operator dalam ekspresi Boolean berkorespondensi langsung dengan satu gerbang logika, dan struktur hierarkis ekspresi menentukan interkoneksi antar gerbang. Proses konversi ini bersifat mekanistik dan dapat dilakukan secara algoritmik \cite{rosen2019}.

\subsection{Prosedur Konversi Sistematis}

Berikut prosedur umum untuk mengkonversi ekspresi Boolean menjadi rangkaian gerbang:

\begin{enumerate}
  \item \textbf{Identifikasi operator terluar}: Operator dengan prioritas paling rendah yang belum dibatasi tanda kurung merupakan gerbang pada tahap keluaran akhir.
  \item \textbf{Dekomposisi rekursif}: Setiap operand dari operator terluar dianalisis secara rekursif. Jika operand merupakan ekspresi komposit, proses diulang. Jika operand merupakan variabel tunggal atau komplemennya, ia menjadi sinyal masukan.
  \item \textbf{Penanganan komplemen}: Setiap komplemen variabel memerlukan satu gerbang NOT pada jalur sinyal masukan.
  \item \textbf{Penanganan percabangan}: Jika satu variabel digunakan di beberapa tempat dalam ekspresi, jalur sinyalnya dicabangkan (\textit{fan-out}).
\end{enumerate}

\subsection{Contoh Konversi Langkah demi Langkah}

\begin{contoh}
\textbf{Konversi $F = A'B + AB'$ (Fungsi XOR)}

\textbf{Langkah 1:} Operator terluar adalah $+$ (OR), sehingga gerbang terakhir (output) adalah OR.

\textbf{Langkah 2:} Operand kiri $A'B$ memerlukan gerbang AND dengan masukan $A'$ dan $B$. Operand kanan $AB'$ memerlukan gerbang AND dengan masukan $A$ dan $B'$.

\textbf{Langkah 3:} $A'$ memerlukan NOT pada $A$, dan $B'$ memerlukan NOT pada $B$.

\textbf{Langkah 4:} Variabel $A$ dicabangkan ke NOT (untuk $A'$) dan ke AND kedua. Variabel $B$ dicabangkan ke AND pertama dan ke NOT (untuk $B'$).

\begin{center}
\begin{tikzpicture}[circuit logic US, huge circuit symbols,
                    every circuit symbol/.style={fill=white, draw=black, thick}]

  % Gates
  \node[not gate] (notA) at (1, 1) {};
  \node[not gate] (notB) at (1, -1.5) {};
  \node[and gate] (and1) at (3, 1.5) {};
  \node[and gate] (and2) at (3, -1) {};
  \node[or gate]  (or1)  at (6, 0.25) {};

  % Inputs
  \node (A) at (-1, 2) {$A$};
  \node (B) at (-1, -1) {$B$};

  % A branching
  \draw (A) -- (0,2) |- (notA.input);
  \draw (0,2) |- (and2.input 1);
  \filldraw (0,2) circle (2pt);

  % B branching
  \draw (B) -- (0,-1) |- (notB.input);
  \draw (0,-1) |- (and1.input 2);
  \filldraw (0,-1) circle (2pt);

  % Inner connections
  \draw (notA.output) |- (and1.input 1);
  \draw (notB.output) |- (and2.input 2);

  \draw (and1.output) -- ++(1,0) |- (or1.input 1);
  \draw (and2.output) -- ++(1,0) |- (or1.input 2);

  \draw (or1.output) -- ++(0.5,0) node[right] {$F = A'B + AB'$};

\end{tikzpicture}
\captionsetup{hypcap=false}
\captionof{figure}{Rangkaian Gerbang untuk $F = A'B + AB'$ (Ekuivalen XOR)}
\end{center}

\textbf{Penghitungan gerbang}: 2 NOT + 2 AND + 1 OR = \textbf{5 gerbang}.
\end{contoh}

\begin{contoh}
\textbf{Konversi $G = (A + B)(C' + D)$}

\textbf{Analisis:} Operator terluar adalah perkalian ($\cdot$) antara dua suku penjumlahan.

\begin{center}
\begin{tikzpicture}[circuit logic US, huge circuit symbols,
                    every circuit symbol/.style={fill=white, draw=black, thick}]

  % Gates
  \node[or gate] (or1) at (2, 1.5) {};
  \node[not gate] (notC) at (0, -0.5) {};
  \node[or gate] (or2) at (2, -1) {};
  \node[and gate] (and1) at (5, 0.25) {};

  % Inputs
  \draw (or1.input 1) -- ++(-0.5,0) node[left] {$A$};
  \draw (or1.input 2) -- ++(-0.5,0) node[left] {$B$};
  \draw (notC.input) -- ++(-0.5,0) node[left] {$C$};
  \draw (or2.input 2) -- ++(-0.5,0) node[left] {$D$};

  % Connections
  \draw (notC.output) |- (or2.input 1);
  \draw (or1.output) -- ++(1,0) |- (and1.input 1);
  \draw (or2.output) -- ++(1,0) |- (and1.input 2);
  \draw (and1.output) -- ++(0.5,0) node[right] {$G$};

\end{tikzpicture}
\captionsetup{hypcap=false}
\captionof{figure}{Rangkaian Gerbang untuk $G = (A + B)(C' + D)$}
\end{center}

\textbf{Penghitungan gerbang}: 1 NOT + 2 OR + 1 AND = \textbf{4 gerbang}.
\end{contoh}

\subsection{Konversi ke Rangkaian NAND-Only}

Dalam praktik industri, rangkaian sering dikonversi ke bentuk yang hanya menggunakan gerbang NAND. Prosedur konversi dari bentuk SOP (AND-OR) ke NAND-NAND menggunakan prinsip involusi ($F = (F')'$) dan hukum De Morgan. Untuk ekspresi dalam bentuk SOP $F = P_1 + P_2 + \ldots + P_k$ (di mana $P_i$ adalah suku perkalian):

\begin{align*}
F &= P_1 + P_2 + \ldots + P_k \\
  &= ((P_1 + P_2 + \ldots + P_k)')' \\
  &= (P_1' \cdot P_2' \cdot \ldots \cdot P_k')' & \text{(De Morgan)}
\end{align*}

\begin{konsep}
Setiap bentuk SOP dua tingkat (AND-OR) dapat dikonversi menjadi rangkaian NAND-NAND dua tingkat yang ekuivalen. Pada tingkat pertama, gerbang AND diganti dengan gerbang NAND. Pada tingkat kedua, gerbang OR diganti dengan gerbang NAND. Prinsip ini dikenal sebagai \logic{konversi NAND-NAND} dan merupakan teknik standar dalam sintesis rangkaian digital.
\end{konsep}

\begin{contoh}
\textbf{Konversi $F = AB + CD$ ke NAND-Only}

\begin{align*}
F &= AB + CD \\
  &= ((AB)' \cdot (CD)' )' & \text{(Involusi + De Morgan)} \\
  &= (AB) \uparrow (CD) & \text{(NAND tingkat 2)}
\end{align*}

Di mana $(AB)$ dan $(CD)$ masing-masing dihasilkan oleh gerbang NAND ditambah inversi. Namun karena NAND pada tingkat 2 sudah menginversi, dan NAND pada tingkat 1 juga menginversi, inversi ganda membatalkan satu sama lain. Hasilnya: gunakan gerbang NAND di semua tingkat --- tingkat 1 menghasilkan $(AB)'$ dan $(CD)'$, dan tingkat 2 menghasilkan $((AB)' \cdot (CD)')' = AB + CD$. Total: \textbf{3 gerbang NAND}.
\end{contoh}

\section{Rangkaian Aritmetika: Half-Adder dan Full-Adder}

Rangkaian aritmetika merupakan salah satu aplikasi fundamental gerbang logika. Operasi penjumlahan biner menjadi dasar bagi seluruh operasi aritmetika dalam prosesor digital, termasuk pengurangan (melalui komplemen dua), perkalian (melalui penjumlahan berulang), dan pembagian. Pemahaman tentang rangkaian penjumlah (\textit{adder}) menghubungkan konsep aljabar Boolean dengan implementasi perangkat keras \cite{rosen2019}.

\subsection{Half-Adder}

\begin{definisi}[Half-Adder]
\logic{Half-adder} adalah rangkaian kombinasional yang menjumlahkan dua bit masukan ($A$ dan $B$) dan menghasilkan dua bit keluaran: \textbf{Sum} ($S$) yang merupakan bit hasil penjumlahan, dan \textbf{Carry} ($C_{out}$) yang merupakan bit bawaan (simpanan) ke posisi bit berikutnya.
\end{definisi}

Persamaan Boolean untuk half-adder diperoleh langsung dari definisi penjumlahan biner:

\begin{align*}
S &= A \oplus B = A'B + AB' \quad \text{(Sum: gerbang XOR)} \\
C_{out} &= A \cdot B \quad \text{(Carry: gerbang AND)}
\end{align*}

Tabel kebenaran half-adder:

\begin{center}
\renewcommand{\arraystretch}{1.2}
\begin{tabular}{|c|c||c|c|l|}
\hline
$A$ & $B$ & $S$ & $C_{out}$ & \textbf{Keterangan} \\ \hline
0 & 0 & 0 & 0 & $0 + 0 = 0$ \\ \hline
0 & 1 & 1 & 0 & $0 + 1 = 1$ \\ \hline
1 & 0 & 1 & 0 & $1 + 0 = 1$ \\ \hline
1 & 1 & 0 & 1 & $1 + 1 = 10_2$ (carry 1) \\ \hline
\end{tabular}
\end{center}

\begin{figure}[H]
\centering
\begin{tikzpicture}[circuit logic US, huge circuit symbols,
                    every circuit symbol/.style={fill=white, draw=black, thick}]
  % XOR for Sum
  \node[xor gate] (xor1) at (2, 1) {};
  \draw (xor1.input 1) -- ++(-0.5,0) coordinate (ja);
  \draw (xor1.input 2) -- ++(-0.5,0) coordinate (jb);
  \draw (xor1.output) -- ++(0.5,0) node[right] {$S = A \oplus B$};

  % AND for Carry
  \node[and gate] (and1) at (2, -1) {};
  \draw (and1.output) -- ++(0.5,0) node[right] {$C_{out} = A \cdot B$};

  % Inputs with branch
  \draw (ja) -- ++(-1,0) node[left] {$A$} coordinate (a);
  \draw (jb) -- ++(-0.5,0) |- (and1.input 1);
  \node (B) at (-0.5,-1.2) {$B$};
  \draw (B) -- ++(0.5,0) coordinate (b) |- (and1.input 2);
  \draw (b) |- (jb);
  \draw (a) ++(0.5,0) |- (and1.input 1);
\end{tikzpicture}
\caption{Diagram Rangkaian Half-Adder}
\label{fig:half-adder}
\end{figure}

Half-adder disebut ``half'' karena hanya mampu menjumlahkan dua bit tanpa memperhitungkan carry dari posisi bit sebelumnya. Keterbatasan ini menjadikan half-adder tidak cukup untuk operasi penjumlahan multi-bit yang memerlukan propagasi carry antar posisi bit.

\subsection{Full-Adder}

\begin{definisi}[Full-Adder]
\logic{Full-adder} adalah rangkaian kombinasional yang menjumlahkan \textbf{tiga} bit masukan: dua operand ($A$ dan $B$) serta satu carry masukan ($C_{in}$) dari posisi bit sebelumnya. Keluarannya adalah Sum ($S$) dan Carry out ($C_{out}$).
\end{definisi}

Persamaan Boolean untuk full-adder:

\begin{align*}
S &= A \oplus B \oplus C_{in} \\
C_{out} &= AB + C_{in}(A \oplus B) = AB + AC_{in} + BC_{in}
\end{align*}

Tabel kebenaran full-adder memiliki $2^3 = 8$ baris:

\begin{center}
\renewcommand{\arraystretch}{1.2}
\begin{tabular}{|c|c|c||c|c|}
\hline
$A$ & $B$ & $C_{in}$ & $S$ & $C_{out}$ \\ \hline
0 & 0 & 0 & 0 & 0 \\ \hline
0 & 0 & 1 & 1 & 0 \\ \hline
0 & 1 & 0 & 1 & 0 \\ \hline
0 & 1 & 1 & 0 & 1 \\ \hline
1 & 0 & 0 & 1 & 0 \\ \hline
1 & 0 & 1 & 0 & 1 \\ \hline
1 & 1 & 0 & 0 & 1 \\ \hline
1 & 1 & 1 & 1 & 1 \\ \hline
\end{tabular}
\end{center}

\begin{konsep}
Full-adder dapat diimplementasikan menggunakan dua half-adder dan satu gerbang OR. Half-adder pertama menjumlahkan $A$ dan $B$, menghasilkan $S_1 = A \oplus B$ dan $C_1 = AB$. Half-adder kedua menjumlahkan $S_1$ dan $C_{in}$, menghasilkan $S = S_1 \oplus C_{in}$ dan $C_2 = S_1 \cdot C_{in}$. Carry akhir: $C_{out} = C_1 + C_2 = AB + (A \oplus B)C_{in}$.
\end{konsep}

\subsection{Ripple-Carry Adder}

Untuk menjumlahkan bilangan biner $n$-bit, $n$ buah full-adder dirangkai secara kaskade (\textit{cascade}), di mana carry output dari setiap posisi bit dihubungkan ke carry input pada posisi bit berikutnya. Rangkaian ini disebut \logic{ripple-carry adder} karena sinyal carry ``mengalir'' (\textit{ripple}) dari posisi bit terendah (LSB) ke posisi bit tertinggi (MSB) \cite{wiki_adder}.

Kelemahan utama ripple-carry adder adalah waktu propagasi yang proporsional terhadap jumlah bit, karena setiap full-adder harus menunggu carry dari full-adder sebelumnya sebelum dapat menghasilkan output yang benar. Untuk penjumlah $n$-bit, waktu propagasi terburuk bergantung pada implementasi (XOR vs NAND, struktur jalur kritis); sebagai aproksimasi sering digunakan orde $O(n)$ atau sekitar $2n$--$3n$ kali delay satu gerbang. Arsitektur yang lebih cepat seperti \textit{carry-lookahead adder} mengatasi keterbatasan ini, tetapi pembahasannya berada di luar cakupan buku ini.

\section[Rangkaian Kombinasional]{\raggedright Rangkaian Kombinasional: Multiplexer dan Decoder}

Selain rangkaian aritmetika, terdapat dua kelas rangkaian kombinasional yang penting dalam arsitektur komputer: \logic{multiplexer} (MUX) dan \logic{decoder}. Kedua rangkaian ini merupakan blok bangunan (\textit{building block}) standar yang digunakan secara masif dalam desain prosesor, memori, dan sistem bus digital \cite{rosen2019}.

\subsection{Multiplexer (MUX)}

\begin{definisi}[Multiplexer]
\logic{Multiplexer} (MUX) adalah rangkaian kombinasional yang memilih satu dari $2^n$ jalur masukan data dan mengarahkannya ke satu jalur keluaran, berdasarkan kombinasi $n$ sinyal pemilih (\textit{select lines}). Multiplexer $2^n$-ke-$1$ memiliki $2^n$ masukan data, $n$ sinyal pemilih, dan 1 keluaran.
\end{definisi}

Multiplexer 2-ke-1 (MUX 2:1) adalah bentuk paling sederhana, dengan 2 masukan data ($D_0, D_1$), 1 sinyal pemilih ($S$), dan 1 keluaran ($Y$).

Ekspresi Boolean MUX 2:1:
\[ Y = S' \cdot D_0 + S \cdot D_1 \]

Ketika $S = 0$, keluaran $Y = D_0$; ketika $S = 1$, keluaran $Y = D_1$. Tabel kebenaran:

\begin{center}
\renewcommand{\arraystretch}{1.2}
\begin{tabular}{|c||c|l|}
\hline
$S$ & $Y$ & \textbf{Keterangan} \\ \hline
0 & $D_0$ & Memilih masukan $D_0$ \\ \hline
1 & $D_1$ & Memilih masukan $D_1$ \\ \hline
\end{tabular}
\end{center}

\begin{contoh}
\textbf{Multiplexer 4-ke-1}

MUX 4:1 memiliki 4 masukan data ($D_0, D_1, D_2, D_3$), 2 sinyal pemilih ($S_1, S_0$), dan 1 keluaran ($Y$).

Ekspresi Boolean:
\[ Y = S_1'S_0' \cdot D_0 + S_1'S_0 \cdot D_1 + S_1 S_0' \cdot D_2 + S_1 S_0 \cdot D_3 \]

\begin{center}
\renewcommand{\arraystretch}{1.2}
\begin{tabular}{|c|c||c|}
\hline
$S_1$ & $S_0$ & $Y$ \\ \hline
0 & 0 & $D_0$ \\ \hline
0 & 1 & $D_1$ \\ \hline
1 & 0 & $D_2$ \\ \hline
1 & 1 & $D_3$ \\ \hline
\end{tabular}
\end{center}
\end{contoh}

\begin{konsep}
MUX dapat digunakan untuk mengimplementasikan fungsi Boolean apa pun. Untuk fungsi $n$ variabel, hubungkan variabel-variabel sebagai sinyal pemilih pada MUX $2^n$:1, dan tetapkan setiap masukan data ke nilai fungsi (0 atau 1) sesuai tabel kebenaran. Teknik ini disebut \logic{implementasi MUX} dan sering dipakai dalam prototipe cepat.
\end{konsep}

\subsection{Decoder}

\begin{definisi}[Decoder]
\logic{Decoder} adalah rangkaian kombinasional yang mengubah kode biner $n$-bit menjadi paling banyak $2^n$ jalur keluaran yang unik, di mana tepat satu jalur keluaran aktif (bernilai 1) pada setiap saat berdasarkan kombinasi masukan. Decoder $n$-ke-$2^n$ memiliki $n$ masukan dan $2^n$ keluaran.
\end{definisi}

Decoder 2-ke-4 memiliki 2 masukan ($A, B$) dan 4 keluaran ($D_0, D_1, D_2, D_3$):

\begin{center}
\renewcommand{\arraystretch}{1.2}
\begin{tabular}{|c|c||c|c|c|c|}
\hline
$A$ & $B$ & $D_0$ & $D_1$ & $D_2$ & $D_3$ \\ \hline
0 & 0 & 1 & 0 & 0 & 0 \\ \hline
0 & 1 & 0 & 1 & 0 & 0 \\ \hline
1 & 0 & 0 & 0 & 1 & 0 \\ \hline
1 & 1 & 0 & 0 & 0 & 1 \\ \hline
\end{tabular}
\end{center}

Persamaan Boolean: $D_0 = A'B'$, $D_1 = A'B$, $D_2 = AB'$, $D_3 = AB$.

\begin{catatan}
Perhatikan bahwa setiap keluaran decoder berkorespondensi tepat dengan satu \textbf{minterm}. Decoder $n$-ke-$2^n$ menghasilkan seluruh $2^n$ minterm dari $n$ variabel masukannya. Oleh karena itu, fungsi Boolean apa pun dapat diimplementasikan menggunakan satu decoder dan satu gerbang OR: hubungkan keluaran decoder yang bersesuaian dengan minterm-minterm fungsi ke masukan gerbang OR.
\end{catatan}

\subsection{Hubungan Decoder dan Multiplexer}

Decoder dan multiplexer saling melengkapi dalam arsitektur digital. Decoder mengonversi kode biner menjadi sinyal aktivasi individual (``satu-dari-banyak''), sedangkan multiplexer melakukan operasi sebaliknya: memilih satu dari banyak sinyal berdasarkan kode biner. Dalam sistem komputer, decoder digunakan untuk pemilihan alamat memori (\textit{address decoding}), sementara multiplexer digunakan untuk pemilihan jalur data pada bus \cite{wiki_multiplexer}.

\section{Analisis Rangkaian: dari Diagram ke Ekspresi Boolean}

Selain kemampuan mengonversi ekspresi Boolean menjadi rangkaian gerbang (sintesis), perancangan digital juga memerlukan proses kebalikannya: menganalisis rangkaian gerbang yang sudah ada dan menentukan ekspresi Boolean serta tabel kebenaran yang diimplementasikannya. Proses ini disebut \logic{analisis rangkaian} dan berperan setara dengan sintesis \cite{rosen2019}.

\subsection{Prosedur Analisis Rangkaian}

Prosedur analisis rangkaian gerbang logika dilakukan secara sistematis dengan langkah-langkah berikut:

\begin{enumerate}
  \item \textbf{Beri label pada semua sinyal}: Identifikasi sinyal masukan primer, keluaran akhir, dan berikan label pada setiap sinyal antara (output dari setiap gerbang internal).
  \item \textbf{Telusuri dari masukan ke keluaran}: Mulai dari gerbang yang langsung terhubung ke masukan primer, tentukan ekspresi Boolean untuk setiap sinyal antara.
  \item \textbf{Substitusi berturut-turut}: Substitusikan ekspresi sinyal antara ke dalam ekspresi gerbang berikutnya, secara progresif menuju keluaran akhir.
  \item \textbf{Sederhanakan (opsional)}: Sederhanakan ekspresi akhir menggunakan aljabar Boolean jika diperlukan.
  \item \textbf{Verifikasi}: Buat tabel kebenaran dari ekspresi yang diperoleh dan bandingkan dengan perilaku rangkaian untuk beberapa kombinasi input.
\end{enumerate}

\subsection{Contoh Analisis Rangkaian}

\begin{contoh}
\textbf{Analisis Rangkaian Tiga Tingkat}

Diberikan rangkaian dengan konfigurasi berikut:
\begin{itemize}
    \item Gerbang 1 (NOT): masukan $A$, keluaran $W = A'$
    \item Gerbang 2 (AND): masukan $W$ dan $B$, keluaran $X = A'B$
    \item Gerbang 3 (AND): masukan $A$ dan $B'$ (NOT pada $B$), keluaran $Y = AB'$
    \item Gerbang 4 (OR): masukan $X$ dan $Y$, keluaran $F = X + Y = A'B + AB'$
\end{itemize}

\textbf{Analisis:}
\begin{align*}
W &= A' \\
X &= W \cdot B = A'B \\
Y &= A \cdot B' \\
F &= X + Y = A'B + AB' = A \oplus B
\end{align*}

Kesimpulan: Rangkaian ini mengimplementasikan fungsi XOR ($A \oplus B$).

\textbf{Verifikasi dengan tabel kebenaran:}
\begin{center}
\renewcommand{\arraystretch}{1.2}
\begin{tabular}{|c|c||c|c|c|c||c|}
\hline
$A$ & $B$ & $W = A'$ & $X = A'B$ & $B'$ & $Y = AB'$ & $F = X + Y$ \\ \hline
0 & 0 & 1 & 0 & 1 & 0 & 0 \\ \hline
0 & 1 & 1 & 1 & 0 & 0 & 1 \\ \hline
1 & 0 & 0 & 0 & 1 & 1 & 1 \\ \hline
1 & 1 & 0 & 0 & 0 & 0 & 0 \\ \hline
\end{tabular}
\end{center}

Hasilnya sesuai dengan tabel kebenaran XOR: output 1 jika dan hanya jika kedua masukan berbeda.
\end{contoh}

\begin{contoh}
\textbf{Analisis Rangkaian NAND-Only}

Diberikan rangkaian yang terdiri dari tiga gerbang NAND:
\begin{itemize}
    \item Gerbang NAND 1: masukan $A$ dan $B$, keluaran $P = (AB)'$
    \item Gerbang NAND 2: masukan $C$ dan $D$, keluaran $Q = (CD)'$
    \item Gerbang NAND 3: masukan $P$ dan $Q$, keluaran $F = (PQ)'$
\end{itemize}

\textbf{Analisis:}
\begin{align*}
P &= (AB)' \\
Q &= (CD)' \\
F &= (PQ)' = ((AB)' \cdot (CD)')' \\
  &= ((AB)')' + ((CD)')' & \text{(De Morgan)} \\
  &= AB + CD & \text{(Involusi)}
\end{align*}

Kesimpulan: Rangkaian tiga gerbang NAND ini mengimplementasikan fungsi $F = AB + CD$, yaitu bentuk SOP dua suku. Ini mengonfirmasi bahwa rangkaian NAND-NAND dua tingkat mengimplementasikan fungsi SOP.
\end{contoh}

\subsection{Analisis Timing dan Propagation Delay}

Selain analisis fungsional (ekspresi Boolean), analisis rangkaian juga mencakup analisis waktu (\textit{timing analysis}). Setiap gerbang logika memiliki waktu propagasi (\textit{propagation delay}), yaitu waktu yang diperlukan untuk perubahan pada masukan untuk tercermin pada keluaran. Delay total rangkaian ditentukan oleh jalur terpanjang (\textit{critical path}) dari masukan ke keluaran.

\begin{definisi}[Critical Path]
\logic{Critical path} (jalur kritis) adalah jalur sinyal dari masukan primer ke keluaran akhir yang melewati jumlah gerbang terbanyak. Delay total rangkaian sama dengan jumlah delay gerbang pada jalur kritis, dan ini menentukan frekuensi operasi maksimum rangkaian \cite{electronics_tutorials_gates}.
\end{definisi}

Sebagai contoh, jika setiap gerbang memiliki delay $t_d$ dan rangkaian memiliki 3 tingkat gerbang, maka delay total jalur kritis adalah $3 \cdot t_d$. Rangkaian dengan lebih sedikit tingkat gerbang (misalnya bentuk SOP dua tingkat) memiliki delay yang lebih kecil dibandingkan rangkaian multi-tingkat, meskipun mungkin memerlukan lebih banyak gerbang. Trade-off antara jumlah gerbang (area/biaya) dan delay (kecepatan) merupakan pertimbangan utama dalam optimasi rangkaian digital.

\subsection{Hubungan dengan bab berikutnya}

Bab 9--11 membahas aljabar Boolean dan realisasinya dalam rangkaian digital. Pada bab-bab berikutnya, pembahasan beralih ke metode pembuktian matematis: penerapan penalaran deduktif pada klaim matematis dan sifat algoritma, sebagai kelanjutan kerangka logis yang telah dibentuk melalui materi logika digital.


% ============================================================
% AKTIVITAS PEMBELAJARAN
% ============================================================

\begin{aktivitas}
  \item \textbf{Rangkaian AND-OR-NOT:} Rancang rangkaian gerbang untuk $F(X,Y,Z) = X(Y' + Z)$ menggunakan gerbang AND, OR, dan NOT. Hitung jumlah gerbang.
  \item \textbf{Konversi NAND-only:} Konversikan rangkaian dari aktivitas sebelumnya ke rangkaian NAND-only. Bandingkan jumlah gerbang sebelum dan sesudah.
  \item \textbf{Half-adder:} Implementasikan half-adder menggunakan hanya gerbang NAND. Gambarkan diagram dan verifikasi dengan tabel kebenaran.
  \item \textbf{Analisis rangkaian:} Diberikan rangkaian: NOT pada $A$ (output $W$), AND($W$, $C$) = $X$, OR($B$, $X$) = $Y$, AND($Y$, $D$) = $F$. Tentukan ekspresi Boolean $F$ dan buat tabel kebenaran.
  \item \textbf{Decoder:} Rancang decoder 3-ke-8 dengan gerbang AND dan NOT. Tulis ekspresi Boolean untuk setiap output.
  \item \textbf{Implementasi dengan decoder:} Implementasikan $G(A,B,C) = \sum m(1, 2, 6, 7)$ dengan: (a) gerbang AND-OR-NOT, dan (b) decoder 3-ke-8 dan satu gerbang OR.
  \item \textbf{Multiplexer:} Rancang multiplexer 4-ke-1 dengan gerbang AND, OR, dan NOT. Verifikasi untuk setiap kombinasi sinyal pemilih.
\end{aktivitas}

% ============================================================
% LATIHAN DAN REFLEKSI
% ============================================================

\begin{latihan}
  \item Representasikan $F = W'XY + WXZ + Y'Z'$ ke diagram rangkaian gerbang. Hitung jumlah gerbang (AND, OR, NOT).
  \item Konversikan rangkaian pada soal sebelumnya ke NAND-only. Tunjukkan langkah konversi dan hitung jumlah gerbang NAND.
  \item Rancang full-adder dan buktikan bahwa $C_{out} = AB + C_{in}(A \oplus B)$ ekuivalen dengan $C_{out} = AB + AC_{in} + BC_{in}$.
  \item Diberikan rangkaian NAND-only: NAND($A,B$) = $P$, NAND($P,P$) = $Q$, NAND($C,D$) = $R$, NAND($R,R$) = $S$, NAND($Q,S$) = $F$. Tentukan $F$ dalam bentuk ekspresi Boolean yang disederhanakan.
  \item Jelaskan perbedaan multiplexer dan decoder. Berikan contoh penggunaan masing-masing dalam arsitektur komputer.
  \item Implementasikan $F(A,B) = A \oplus B$ (XOR) dengan: (a) hanya gerbang NAND, (b) hanya gerbang NOR. Gambarkan kedua rangkaian dan bandingkan jumlah gerbang.
  \item \textbf{Refleksi:} (a) Mengapa banyak pabrikan semikonduktor mengonversi rangkaian ke NAND-only? Apa keuntungan dan kerugiannya dibanding campuran AND-OR-NOT? (b) Bagaimana critical path memengaruhi trade-off antara kecepatan dan biaya (jumlah gerbang) dalam desain rangkaian digital?
\end{latihan}

% ============================================================
% ASESMEN
% ============================================================

\begin{asesmen}
\textbf{Instrumen Evaluasi Kinerja (Sub-CPMK 3.3)}

\textbf{A. Pilihan Ganda}
\begin{enumerate}
  \item Gerbang yang output-nya 1 jika dan hanya jika jumlah masukan bernilai 1 ganjil adalah:
  \begin{enumerate}
    \item OR
    \item AND
    \item XOR (Exclusive-OR)
    \item NOR
  \end{enumerate}
  Jawaban: (c)

  \item Gerbang yang secara individual lengkap secara fungsional (\textit{universal gate}) adalah:
  \begin{enumerate}
    \item AND dan OR
    \item XOR dan XNOR
    \item NAND dan NOR
    \item AND dan NOT
  \end{enumerate}
  Jawaban: (c)

  \item Rangkaian half-adder terdiri dari:
  \begin{enumerate}
    \item 1 gerbang AND dan 1 gerbang OR
    \item 1 gerbang XOR dan 1 gerbang AND
    \item 2 gerbang AND dan 1 gerbang OR
    \item 1 gerbang NAND dan 1 gerbang NOR
  \end{enumerate}
  Jawaban: (b)

  \item Decoder $n$-ke-$2^n$ menghasilkan:
  \begin{enumerate}
    \item $n$ sinyal output
    \item $2n$ sinyal output
    \item $2^n$ sinyal output, masing-masing berkorespondensi dengan satu minterm
    \item $n^2$ sinyal output
  \end{enumerate}
  Jawaban: (c)

  \item Jika rangkaian memiliki 4 tingkat gerbang dan setiap gerbang delay $t_d = 5$ ns, delay total jalur kritis adalah:
  \begin{enumerate}
    \item 5 ns
    \item 10 ns
    \item 15 ns
    \item 20 ns
  \end{enumerate}
  Jawaban: (d)
\end{enumerate}

\textbf{B. Tugas Praktik}
\begin{enumerate}
  \item Buktikan gerbang NOR lengkap secara fungsional dengan mengimplementasikan NOT, AND, dan OR hanya menggunakan gerbang NOR. Gambarkan rangkaian masing-masing.
  \item Rancang full-adder 1-bit dan rangkaikan empat full-adder menjadi penjumlah 4-bit (ripple-carry adder). Jelaskan propagasi sinyal carry.
  \item Diberikan rangkaian campuran gerbang AND, OR, NOT, dan NAND. Analisis rangkaian: tentukan ekspresi Boolean output, sederhanakan, dan verifikasi dengan tabel kebenaran.
\end{enumerate}

\textbf{Rubrik}: Merujuk Lampiran Buku.
\end{asesmen}

% ============================================================
% CHECKLIST KOMPETENSI
% ============================================================

\begin{checklist}
  \item Saya memahami fungsi dan simbol gerbang logika dasar (AND, OR, NOT).
  \item Saya mengenal fungsi dan simbol gerbang turunan (NAND, NOR, XOR, XNOR) beserta tabel kebenarannya.
  \item Saya memahami kelengkapan fungsional dan dapat membuktikan universalitas gerbang NAND dan NOR.
  \item Saya mampu mengonversi ekspresi Boolean SOP ke rangkaian gerbang secara sistematis.
  \item Saya mampu mengonversi rangkaian AND-OR-NOT ke rangkaian NAND-only atau NOR-only.
  \item Saya memahami struktur dan persamaan Boolean half-adder dan full-adder.
  \item Saya mengenal multiplexer dan decoder serta hubungannya dengan minterm.
  \item Saya mampu menganalisis rangkaian gerbang dan menentukan ekspresi Boolean yang diimplementasikan.
  \item Saya memahami critical path dan pengaruhnya terhadap kecepatan rangkaian.
\end{checklist}

% ============================================================
% RANGKUMAN
% ============================================================

\begin{rangkuman}
Bab ini membahas realisasi fisik konsep aljabar Boolean dalam bentuk rangkaian gerbang logika digital, menjembatani teori matematika dengan implementasi perangkat keras.

\textbf{Poin Kunci:}
\begin{itemize}
  \item Tiga gerbang dasar (AND, OR, NOT) mengimplementasikan tiga operasi fundamental aljabar Boolean. Empat gerbang turunan (NAND, NOR, XOR, XNOR) memberikan efisiensi tambahan dalam implementasi.
  \item Gerbang NAND dan NOR masing-masing bersifat \textit{lengkap secara fungsional} (universal): seluruh fungsi Boolean dapat diimplementasikan hanya menggunakan satu jenis gerbang. Hal ini memiliki implikasi penting dalam efisiensi fabrikasi semikonduktor.
  \item Konversi ekspresi Boolean ke rangkaian gerbang bersifat mekanistik --- setiap operator menjadi satu gerbang, dan struktur ekspresi menentukan interkoneksi. Rangkaian SOP (AND-OR) dapat dikonversi ke NAND-NAND, dan POS (OR-AND) ke NOR-NOR.
  \item Rangkaian aritmetika (half-adder, full-adder, ripple-carry adder) merupakan blok bangunan dasar untuk operasi penjumlahan biner dalam prosesor.
  \item Multiplexer memilih satu dari banyak sinyal masukan berdasarkan kode pemilih, sementara decoder mengonversi kode biner menjadi sinyal aktivasi individual (satu-dari-banyak). Keduanya dapat digunakan untuk mengimplementasikan fungsi Boolean.
  \item Analisis rangkaian (dari diagram ke ekspresi) dan analisis timing (critical path) merupakan keterampilan penting dalam verifikasi dan optimasi desain digital.
\end{itemize}

\textbf{Kata Kunci}: gerbang logika, AND, OR, NOT, NAND, NOR, XOR, XNOR, gerbang universal, kelengkapan fungsional, half-adder, full-adder, ripple-carry adder, multiplexer, decoder, analisis rangkaian, critical path
\end{rangkuman}

\ifSubfilesClassLoaded{
  \renewcommand{\bibname}{Daftar Pustaka}
  \bibliographystyle{plain}
  \bibliography{../references}
}{}
\end{document}
